\section{The general \texorpdfstring{$\chi$}{chi}-twisted descent}
\label{sec:general}

The proof of Theorem~\ref{thm:brow} in Section~\ref{sec:weight3} is a two-page Kummer
ledger. It turns out that the ledger is an artefact of Ap\'ery's own choice of weight: if
one replaces the nested alternating sum $\sum_{m\le k}u(n,m)$ by a plain harmonic monomial
--- which Section~\ref{sec:minimal} shows is possible, and in a startlingly short form ---
the whole argument collapses to three lines, and generalises at once to all weights, to
several sequences, and to a twist by a quadratic character.

\subsection{Lemma K: Frobenius descent of a harmonic letter}

\begin{lemma}[Lemma K]\label{lem:K}
\Proved\ Let $\chi$ be a completely multiplicative arithmetic function (a Dirichlet
character; $\chi\equiv1$ is allowed), let $r\ge1$ and $y\ge0$, and set
$\Kchi ry=\sum_{m=1}^{y}\chi(m)m^{-r}$. Then, as an identity in $\Q$,
\[
   p^{r}\,\Kchi ry\;=\;\chi(p)\,\Kchi r{\lfloor y/p\rfloor}\;+\;p^{r}\,T,
   \qquad
   T:=\sum_{\substack{m\le y\\ p\nmid m}}\frac{\chi(m)}{m^{r}}\;\in\;\Zp .
\]
\end{lemma}

\begin{proof}
Split the sum according to $p\mid m$. The $p$-divisible part is
\[
   \sum_{j\le\lfloor y/p\rfloor}\frac{\chi(jp)}{(jp)^{r}}
   \;=\;\chi(p)\,p^{-r}\sum_{j\le\lfloor y/p\rfloor}\frac{\chi(j)}{j^{r}}
   \;=\;\chi(p)\,p^{-r}\,\Kchi r{\lfloor y/p\rfloor},
\]
using complete multiplicativity. The remaining terms have $p$-integral denominators, so
their sum is $T\in\Zp$. Multiply by $p^r$.
\end{proof}

\emph{This single line is the entire source of the character in Conjecture~\ref{conj:master}.}
Sequences whose second solution needs $\chi$-twisted harmonic letters acquire $\chi(p)$;
sequences whose second solution is a pure $\zeta$-harmonic monomial acquire $+1$.

\begin{corollary}[Corollary K2]\label{cor:K2}
\Proved\ Let $M=\prod_{t=1}^{s}K^{(r_t)}_{\chi_t}(x_t)$ with $\sum_tr_t=w$ and each
$K^{(r_t)}_{\chi_t}(\lfloor x_t/p\rfloor)\in\Zp$. Then
\[
   p^{w}M\;\equiv\;\Bigl(\prod_t\chi_t(p)\Bigr)\prod_tK^{(r_t)}_{\chi_t}
   \bigl(\lfloor x_t/p\rfloor\bigr)\pmod p .
\]
\end{corollary}

\begin{proof}
$p^wM=\prod_t\bigl(p^{r_t}K^{(r_t)}_{\chi_t}(x_t)\bigr)
=\prod_t\bigl(\chi_t(p)K^{(r_t)}_{\chi_t}(\lfloor x_t/p\rfloor)+p^{r_t}T_t\bigr)$
by Lemma~\ref{lem:K}. Expanding, every term other than the pure product carries a factor
$p$, and every factor is $p$-integral by hypothesis.
\end{proof}

\subsection{Theorem LB}

\begin{definition}[admissible pair]\label{def:admissible}
Let $p\ge5$. An \emph{admissible pair} consists of
\begin{itemize}[leftmargin=1.4em]
\item a summand $S(n,k)=\prod_i\bin{L_i(n,k)}{M_i(n,k)}$ with $L_i,M_i$ integer linear
  forms, and $A(n)=\sum_kS(n,k)$;
\item a weight $w(n,k)=\sum_jc_j\prod_tK^{(r_{jt})}_{\chi}(x_{jt})$, a $\Q$-combination of
  monomials of total degree $\sum_tr_{jt}=w$ in letters attached to one fixed quadratic
  character $\chi$ and to the trivial character, with the $x_{jt}$ integer linear forms in
  $(n,k)$;
\item the second row $B(n)=\sum_kS(n,k)w(n,k)$.
\end{itemize}
\end{definition}

Fix $1\le a<p$ and $0\le r,s<p$, and write $n=ap+r$, $k=bp+s$. The hypotheses are:

\begin{description}[leftmargin=3.0em,style=sameline]
\item[(H1)] \emph{Lucas/carry dichotomy.} For every $k$, either $p\mid S(n,k)$, or
  $S(n,k)\equiv S(a,b)S(r,s)\pmod p$.
\item[(H2)] \emph{Product region.} The surviving set is $\{0\le b\le a\}\times\Sigma_r$
  for some $\Sigma_r\subseteq[0,p)$, and $\sum_{s\in\Sigma_r}S(r,s)\equiv A(r)\pmod p$.
\item[(H3)] \emph{Digit compatibility.} On the surviving set,
  $\lfloor x(n,k)/p\rfloor=x(a,b)$ for every argument form $x$ occurring in $w$.
\item[(H4)] \emph{Tameness.} Every argument form satisfies $0\le x(n,k)\le n$ on the
  summation range, and $c_j\in\Zp$ for all $j$.
\item[(H5)] \emph{$\chi$-homogeneity.} Every monomial of $w$ contains exactly the same
  number $e\in\{0,1\}$ of $\chi$-letters, all other letters carrying the trivial
  character.
\end{description}

\begin{theorem}[Theorem LB]\label{thm:LB}
\Proved\ Let $p\ge5$ and let $(S,w,A,B)$ be an admissible pair satisfying
\textup{(H1)--(H5)}. Then for all $1\le a<p$ and $0\le r<p$, with $n=ap+r$,
\[
   p^{w}\,B(ap+r)\;\equiv\;\chi(p)^{e}\,B(a)\,A(r)\pmod p .
\]
\end{theorem}

\begin{proof}
By (H4), every argument satisfies $x(n,k)\le n<p^2$, hence $\lfloor x/p\rfloor\le a<p$ and
every letter $K^{(r)}(\lfloor x/p\rfloor)$ is $p$-integral; likewise $x(a,b)\le a<p$, so
%% FIXED [MINOR]: the one-line reason for $p^w w(n,k) \in \Zp$ was missing.
$w(a,b)\in\Zp$. For $p^ww(n,k)$: each letter $K^{(r)}_\chi(x)$ with $x\le n<p^2$ has
$\vp\ge-r$ by Lemma~\ref{lem:K}, a monomial $\prod_tK^{(r_{jt})}_\chi$ of total degree $w$
therefore has $\vp\ge-w$, and (H4) puts the coefficients $c_j$ in $\Zp$; hence
$p^ww(n,k)\in\Zp$ for \emph{every} $k$.
Corollary~\ref{cor:K2} then gives, for every $k$,
\begin{equation}\label{eq:LB-letter}
   p^{w}w(n,k)\;\equiv\;\chi(p)^{e}\,\widetilde w(n,k)\pmod p,
   \qquad
%% FIXED [MINOR]: the letter is the chi-twisted $K^{(r)}_\chi$; the subscript was dropped
%% FIXED (log typo at LBW_GENERAL.md:271).
   \widetilde w(n,k):=\sum_jc_j\prod_tK^{(r_{jt})}_{\chi}\bigl(\lfloor x_{jt}/p\rfloor\bigr),
\end{equation}
and (H3) upgrades $\widetilde w(n,k)=w(a,b)$ on the surviving set. Now split
\[
   p^{w}B(n)\;=\;\sum_{k\,:\,p\mid S(n,k)}S(n,k)\,p^{w}w(n,k)
   \;+\;\sum_{k\text{ surviving}}S(n,k)\,p^{w}w(n,k).
\]
The first sum is $\equiv0\pmod p$, because each $p^ww(n,k)$ is $p$-integral and each
$S(n,k)$ is divisible by $p$. In the second, substitute (H1) and
\eqref{eq:LB-letter}:
\[
   p^{w}B(n)\;\equiv\;\chi(p)^{e}\sum_{b=0}^{a}\sum_{s\in\Sigma_r}S(a,b)S(r,s)w(a,b)
   \;=\;\chi(p)^{e}\Bigl(\sum_{b=0}^{a}S(a,b)w(a,b)\Bigr)\Bigl(\sum_{s\in\Sigma_r}S(r,s)\Bigr),
\]
which by (H2) is $\equiv\chi(p)^{e}B(a)A(r)\pmod p$.
\end{proof}

\subsection{Remarks on the hypotheses}

\begin{remark}[three lines instead of two pages]\label{rem:threeLines}
With $e=0$, $S(n,k)=A(n,k)$ and the weight $w(n,k)=2\HH3n-\HH3k$ of
Theorem~\ref{thm:minimal} below, Theorem~\ref{thm:LB} \emph{reproves
Theorem~\ref{thm:brow}}. Hypotheses (H1)--(H3) for this instance are exactly
Lemmas~\ref{lem:summand} and \ref{lem:fibre}; (H4) holds because the arguments are $n$ and
$k$, both $\le n$, and the coefficients are the integers $2,-1$; (H5) is vacuous. No
Lemma~\ref{lem:V}, no Corollary~\ref{cor:tailfact}, no Kummer ledger. The entire content
of the long proof is absorbed into the choice of weight.
\end{remark}

\begin{remark}[what the mechanism actually needs]\label{rem:needs}
In the language of Section~\ref{sec:weight3}: ``summand with Lucas factorisation and carry
annihilation'' is (H1)$+$(H2), which is available for all fifteen sporadic sequences
because it is Malik--Straub's theorem \cite{MalikStraub}; ``harmonic weight of depth one''
is the shape of $w$; and the per-carry valuation bound $\vp\ge2$ of
Remark~\ref{rem:square} is \emph{not} needed --- it is replaced by (H4), tameness, which
makes $p^ww$ $p$-integral outright. Where tameness fails --- because an argument form
reaches $2n$, so that $p^w w(n,k)$ acquires a pole on the vanishing layer --- one needs a
substitute valuation bound; see Section~\ref{ssec:lemmaD}.
\end{remark}

\begin{remark}[scope: single digit, modulo $p$]\label{rem:LBscope}
Theorem~\ref{thm:LB} gives the modulo-$p$, single-digit statement. The master form of
Conjecture~\ref{conj:master} --- all $n$, modulo $p^w$ --- is \Verified\ but not proved
here. The gap is exactly the one left open at weight three
(Remark~\ref{rem:not-proved}).
\end{remark}

\begin{corollary}[integrality]\label{cor:integrality}
Iterating the ratio form of Conjecture~\ref{conj:master} over the base-$p$ digits gives
$\vp(B_n)\ge-w\lfloor\log_pn\rfloor$, i.e.\ $d_n^{w}B_n\in\Zp$ --- the coefficient-level
shadow of $p$-integrality of the mirror map, see Section~\ref{sec:discussion}. Independent
of the conjecture, this bound is \VerifiedR{15 sequences $\times$ 16 primes $\times$
$n\le400$, zero failures}.
\end{corollary}

\begin{remark}[the formalised hypotheses differ, and are weaker]\label{rem:leanH}
The Lean~4 formalisation of Theorem~\ref{thm:LB} (Section~\ref{sec:lean}) replaces
(H1)$+$(H2) by the single unconditional congruence
$S(ap+r,bp+s)\equiv S(a,b)S(r,s)\pmod p$ for \emph{all} $a,b$ together with the vanishing
$S(n,k)=0$ for $k>n$. In the carrying regime both sides vanish modulo $p$, so the
dichotomy is absorbed and no product-region argument occurs anywhere in the machine proof;
the double sum factors because a \emph{full} block $\{0\le k<(a+1)p\}$ splits, and the
extension from $\{0\le k\le n\}$ to the full block is an equality of integers rather than
a congruence. It also replaces the counting form of (H5) by the equivalent but more
flexible requirement that the product of the characters of the letters of each monomial,
evaluated at $p$, be constant.
\end{remark}
